\section{Study Design}
\label{sec:design}

The study evaluates the concept of Section~\ref{sec:concept} on published data, without new model calls.
It has three parts: RQ1 on crisis labels, a descriptive analysis of expert uncertainty in CounselBench, and RQ2 on
reply verdicts; a fourth part combines RQ1 and RQ2 in a demonstration of the escalation decoder.
This section describes the data, the alignment rules, the pre-registration, the compared systems, the metrics and the
measures taken for reproducibility.

\subsection{Datasets and Licences}
\label{sec:design-data}

Table~\ref{tab:datasets} summarises the three datasets.
All of them are used exactly as published; derived tables are produced only by scripts and are never edited by hand
(decision D2 in Appendix~\ref{app:decisions}).

\begin{table}[t]
\caption{Datasets analysed in this paper. Units are the items that enter the analyses after the alignment rules of
Section~\ref{sec:design-alignment}.}
\label{tab:datasets}
% Generated by paper/scripts/export_tables.py from results/. Do not edit by hand.
\begingroup\footnotesize
\begin{tabular}{p{0.17\linewidth}p{0.14\linewidth}rp{0.27\linewidth}p{0.2\linewidth}}
\toprule
Dataset & Unit & Units & Raters & Licence \\
\midrule
BHH crisis labels & user message & \RoneNItemsAll & four human annotators; three LLMs $\times$ three runs & none stated (research use with citation) \\
BHH appropriateness & reply to a message & \RtwoNItemsAll & two human experts; three LLM judges $\times$ three runs & none stated (research use with citation) \\
CounselBench-Eval & reply to a question & \CbNRepliesAll & five clinical experts per reply & CC BY-NC-ND 4.0 \\
\bottomrule
\end{tabular}
\endgroup

\end{table}

\paragraph{BHH crisis labels.}
The BHH validation set contains \RoneNItemsAll{} user messages \citep{arnaizrodriguez2026help}.
Each message was labelled by four human annotators and, in the published files, by three LLMs in three independent runs
each.
Labels use six crisis categories (suicidal ideation, self-harm, anxiety crisis, violent thoughts, substance use or
withdrawal, and risk-taking behaviours) and a no-crisis category.
In the published files, \RoneNLlmMissingLabelsAll{} LLM labels across \RoneNItemsWithMissingLabels{} messages are empty;
they stay missing and are never read as ``no crisis'' (D13), so between \RoneMinLlmObserved{} and
\RoneMaxLlmObserved{} LLM labels are observed per message.

\paragraph{BHH appropriateness.}
For the same messages, BHH publishes replies that are rated on the five-point appropriateness scale by two human
experts and by three LLM judges (gpt-4o-mini, gpt-5-nano and a Llama model), each in three runs.
The items were sampled towards the lowest judge scores, so unsafe replies are over-represented relative to ordinary
use.
After alignment, \RtwoNItemsAll{} items have scores from both experts and all three judges; one item has only two runs
of one judge, and this is kept rather than padded.
Of these items, \RtwoNHumanUnsafeAll{} are \emph{human-unsafe}, meaning that at least one expert rated the reply
\word{Inappropriate} or worse.

\paragraph{CounselBench-Eval.}
CounselBench-Eval contains \CbNRepliesAll{} replies to patient questions, written by three LLMs and by online human
therapists, each rated by five clinical experts \citep{li2026counselbench}.
The ratings cover four Likert dimensions (overall quality, empathy, specificity and toxicity), a judgement of medical
advice (yes, no or not sure), factual consistency, and free-text copies of sentences that an expert considered risky.
This part is descriptive only: no LLM judge is evaluated on it (D6).

\paragraph{Licences and ethics of reporting.}
CounselBench is licensed CC BY-NC-ND 4.0, and BHH publishes its data without a licence, which we treat as permission for
research use with citation only (D7).
Neither raw nor derived data are redistributed, and the paper never quotes or paraphrases a message or a reply.
All results consist of item identifiers and numbers; worked cases are described by identifiers, degrees, words and
actions.

\subsection{Use as Published and Alignment Rules}
\label{sec:design-alignment}

The guiding principle is that the published files are the data.
Cleaning and exclusions happen only in the loaders, follow explicit rules and are logged with counts.
Three rules were needed for BHH; Appendix~\ref{app:decisions} gives the full decision log.

\paragraph{Item alignment (D11).}
BHH stores crisis labels and appropriateness scores in separate files that share no identifier.
An item is identified by its published position, and positional alignment is accepted only after its content has
been verified per item and per file: the message text must match exactly after joining its turns, and the reply must
match after a comparison-only normalisation of white space or be near-identical (similarity of at least $0.99$).
Cells that hold spreadsheet artefacts instead of text are accepted only if the other field of the item verifies.
No item had to be excluded under this rule.

\paragraph{Per-run judge scores (D12).}
The sampled judge files store only an aggregate score per item.
The per-run scores are recovered from each judge's own merged result file by an exact match on message and reply,
requiring exactly one match and an identical aggregate.
Runs are never imputed, duplicated or padded.

\paragraph{Missing labels (D13).}
As stated above, missing LLM crisis labels are excluded from both the numerator and the denominator of the crisis
degree, and the number of observed labels is kept per item.

\subsection{Pre-Registration}
\label{sec:design-prereg}

Both analyses were pre-registered in the sense of \citet{nosek2018preregistration}: the analysis plan was written and
committed to the project repository before any analysis code was written or run.
The plan for RQ2 was frozen in commit \PlanRtwoCommit{} on \PlanRtwoDate{}; the plan for RQ1, the CounselBench analysis
and the escalation demo was frozen in commit \PlanRoneDemoCommit{} on \PlanRoneDemoDate{}.
Each plan fixes the data, the systems, the metrics, the hypotheses with their direction, the strata and the sensitivity
analyses.
The plans were never edited afterwards; every deviation is recorded with its reason in the generated summary of the
respective analysis.

All hypothesis tests are exploratory: there is no correction for multiplicity, and a verdict is not a confirmatory
claim (D15).
A pre-registered result keeps its recorded verdict, even where a later reading would suggest a different rule.

\subsection{Compared Systems}
\label{sec:design-systems}

\paragraph{RQ2: reply verdicts.}
Five systems turn the judge runs of a reply into a word on the appropriateness scale.
\begin{itemize}
  \item \textbf{S1 (single judge)}: the mean of one judge's runs, decoded with the 2-tuple; one system per judge.
  \item \textbf{S2 (crisp jury)}: the mean of all runs of all judges, decoded with the 2-tuple, as in an LLM jury.
  \item \textbf{S3 (fuzzy panel)}: the panel set of all judges (Section~\ref{sec:concept-encoder}), decoded with the
        codebook and the Jaccard similarity.
  \item \textbf{S4 (OWA)}: the OWA of the judges' run means with orness $0.25$, decoded with the 2-tuple.
  \item \textbf{S5 (LWA)}: the LWA of the per-judge panel sets with self-consistency weights, decoded like S3.
\end{itemize}
The human reference of an item is the mean of the two expert scores, decoded with the 2-tuple, together with the
type-1 IAA set of the two scores.

\paragraph{RQ1: crisis labels.}
For each message, the human crisis degree is $\mu_H = k_H/4$, where $k_H$ is the number of human annotators who chose a
crisis category, and the LLM crisis degree is $\mu_L = k_L/n_L$ over the observed LLM labels.
A message is \emph{human-split} if $0 < \mu_H < 1$, that is, if the human annotators did not agree.
The crisp alternative is the majority label of each group.

\paragraph{Escalation demo.}
Three pipelines feed the rule base of Section~\ref{sec:concept-rules}.
\emph{P-fuzzy} uses the crisis-word memberships of $\mu_L$ and the S3 similarities.
\emph{P-crisp} uses the LLM majority label (crisis as \word{Clear}, no crisis as \word{None}, a tie as \word{Possible})
and the S2 word, each with membership one.
The \emph{oracle} applies the same rule base to the human inputs: the memberships of $\mu_H$ and the similarities of
the experts' IAA set.
Only P-fuzzy and P-crisp could be deployed, because they use no human data.
Because the oracle uses the same rule table, the comparison isolates the input representation, not the quality of the
rule table.

\subsection{Metrics and Statistical Analysis}
\label{sec:design-metrics}

\paragraph{RQ2 metrics.}
Agreement with the human word (M1) is measured as exact agreement, agreement within one level and quadratic-weighted
Cohen's kappa \citep{cohen1968weighted}.
Fuzzy agreement (M2) is the Jaccard similarity between the system's set and the human set; crisp systems are
represented by the IAA set of their value.
Safety (M3, the primary metric) is the recall and precision of flagging human-unsafe items, where a system flags an item
if its word is \word{Harmful/Unacceptable} or \word{Inappropriate}; for S3 and S5, variant M3b adds the
escalate-when-unsure rule of Section~\ref{sec:concept-decoder}.
Stability (M4) is the share of items whose word changes when one judge is removed (M4a), and the share of items whose
word changes when each judge contributes one randomly drawn run, averaged over \ResampleDraws{} draws (M4b).
The uncertainty signal (M5) is the Spearman correlation of a system's uncertainty with the disagreement between the two
experts and with the system's absolute error; uncertainty is the standard deviation of all runs for S2, and one minus the
largest Jaccard similarity or the centroid width for S3 and S5.

\paragraph{RQ1, CounselBench and demo metrics.}
RQ1 reports the distributions of $\mu_H$ and $\mu_L$, the mean $\mu_L$ at each level of $\mu_H$, and a category-level
fuzzy agreement, $\sum_c \min(p_H(c), p_L(c)) / \sum_c \max(p_H(c), p_L(c))$, where $p(c)$ is the share of raters who
chose category $c$.
The CounselBench analysis reports how often expert panels split on medical advice, how often their ratings span two or
more levels, how often a decoded toxicity panel is contested, and how often exactly one expert flagged a risky
sentence.
The demo reports agreement of each pipeline with the oracle (exact and quadratic-weighted kappa on the action index) and
the rates of under- and over-escalation, that is, of choosing a less or a more severe action than the oracle.

\paragraph{Hypotheses.}
The plans state six directional hypotheses.
H1a: among human-split messages, $\mu_H$ and $\mu_L$ are positively correlated.
H1b: crisp majority labels disagree more often on human-split messages than on unanimous ones.
H2a: S3 and S5 flag human-unsafe replies with at least the recall of S2, and with a precision at most $0.10$ lower.
H2b: S3 and S5 change their word less often than S2 when one judge is removed.
H2c: the uncertainty of S3 and S5 correlates more strongly with expert disagreement than the run spread of S2.
H3: P-fuzzy under-escalates on oracle A2 or A3 items no more often than P-crisp.

\paragraph{Confidence intervals and verdicts.}
Every estimate is reported with a 95\% percentile confidence interval from an item bootstrap with \BootB{} resamples
and the fixed seed \BootSeed{} \citep{efron1993bootstrap}.
Comparisons between systems use paired differences, in which both systems are evaluated on the same resampled items.
A hypothesis is \emph{supported} if the confidence interval of its stated quantity lies entirely on the hypothesised
side of its boundary, \emph{not supported} if it lies entirely on the other side, and \emph{inconclusive} otherwise; if two
systems behave identically, the interval collapses onto the boundary and the verdict stays inconclusive (D15).
All metrics are also reported for messages labelled as no crisis and as a crisis, and the pre-registered sensitivity
analyses vary the codebook FOU width, the OWA orness and the contested threshold.

\subsection{Reproducibility}
\label{sec:design-repro}

The analyses run offline and with fixed seeds, using only the loaders, the CWW library and one script per part.
Running each script twice produced byte-identical metric files, and an automated test suite checks the library and the
decision rules on synthetic data.
Every number in this paper is a macro that a script exports from the result files, and a further script rejects any
number typed by hand into the text; Appendix~\ref{app:prototype} documents the repository and how to reproduce all
results.
